\begin{abstract}
The program models the universe as one four-dimensional hyperbolic honeycomb, $\{3,4,3,4\}$, whose every site holds
24 ternary values updated by one deterministic, reversible, integer rule, with physics read on a flat
three-dimensional surface at its cusp. This paper states the model as it now stands, in one pass: the base, the rule
piece by piece with the reason for each, what the rule produces against computed controls, and what is still open.
Since \emph{The Missing Amplitude} the rule has gained exact amplitudes in $\mathbb{Z}[\omega][1/2]$, a working
vacuum, a compact $U(1)$ light, and a gravity found by summing one trit per link, which bends light by 2, travels at
$c$, forms horizons at $2GM$, obeys an area law, and whose spin 2 is forced by the unlabeled docks. The central new
result is a no-go and a way out: on the rule every vibe is a lineon, line-only motion contradicts observation by
some 17 orders of magnitude, and a stated rule change (a filled sea, a fermionic dock mixer, a swap coin, and the prime
7) gives every excitation one limiting speed and exact relativistic motion in 3d. Binding under that change reduces,
through a chain of theorems, to one dynamical question about a 120-element gauge register on the husk. A changelog
closes the paper. The suite holds 1,474 experiments, every number here cites one, and every failure is kept. Claude
Code was used throughout, with Claude Opus 5 and Opus 5.5, to write the code and this paper.
\end{abstract}
